\documentclass{ieeeaccess}
\usepackage{cite}
\usepackage{amsmath,amssymb,amsfonts}
\usepackage{graphicx}
\usepackage{booktabs}
\usepackage{multirow}
\usepackage{url}
\usepackage[hidelinks]{hyperref}

\setcounter{dbltopnumber}{2}
\renewcommand{\dbltopfraction}{0.9}
\renewcommand{\dblfloatpagefraction}{0.65}
\makeatletter
\setlength{\@dblfptop}{0pt}
\setlength{\@dblfpbot}{0pt plus 1fil}
\makeatother
\makeatletter
\let\IEEEAccessOriginalThebibliography\thebibliography
\renewcommand{\thebibliography}[1]{%
  \IEEEAccessOriginalThebibliography{#1}%
  \setlength{\parskip}{0pt}%
  \setlength{\itemsep}{0pt}%
  \setlength{\parsep}{0pt}%
}
\makeatother

\newsavebox{\paperfigbox}
\newdimen\xfigwd
\newcommand{\setbenchmarkfigure}[2]{%
  \sbox{\paperfigbox}{\includegraphics[#2]{#1}}%
  \global\xfigwd=\wd\paperfigbox
  \usebox{\paperfigbox}%
}
\newcommand{\architecturecaption}[1]{%
  \parbox[t]{0.80\columnwidth}{\raggedright #1}%
}
\newcommand{\resultcaption}[1]{%
  \parbox[t]{0.84\textwidth}{\raggedright #1}%
}
\newcommand{\benchmarkfigure}[2]{%
  \IfFileExists{results/benchmark_pilot_5seed_5round_20261106/figures/#1}{%
    \setbenchmarkfigure{results/benchmark_pilot_5seed_5round_20261106/figures/#1}{#2}%
  }{%
    \IfFileExists{../results/benchmark_pilot_5seed_5round_20261106/figures/#1}{%
      \setbenchmarkfigure{../results/benchmark_pilot_5seed_5round_20261106/figures/#1}{#2}%
    }{%
      \PackageError{benchmarkfigures}{Required figure #1 not found}{%
        Generate the benchmark figures with experiments.plot_benchmark_results.}%
    }%
  }%
}

\begin{document}

\history{Manuscript prepared November 2026}
\doi{}

\title{Adaptive Contribution Scoring for Heterogeneous Federated Fraud Detection}

% The author list follows the prior manuscript; all authors must verify their
% contribution, affiliation, contact details, and approval of the final draft.
\author{\uppercase{Harsh Pathak}\authorrefmark{1},
\uppercase{Siddhant Rathi}\authorrefmark{1},
\uppercase{Aryan Kumar}\authorrefmark{1}, and
\uppercase{Shaurya Singh Chauhan}\authorrefmark{1}}
\address[1]{Student, VIT-AP University, Amaravati, Andhra Pradesh, India}
\corresp{Corresponding author: Harsh Pathak (pathakharsh2004@gmail.com).}

\begin{abstract}
We study a federated fraud-detection prototype for clients with
different feature sets. Each simulated client keeps its own encoder; only a
shared prediction torso is aggregated. We test an adaptive contribution score
that pairs local validation quality with update reliability, and hard-fraud
novelty with leave-one-client-out complementarity. Its client-specific gate
also uses temporal utility. The same-schema comparison covers FedAvg,
FedProx, trimmed mean, and contribution-aware aggregation; the heterogeneous
comparison uses encoder+FedAvg and encoder+adaptive pairs. We evaluate the
MLG-ULB credit-card dataset and a separate 10,000-row Kaggle dataset
\cite{miadul_dataset}, with five seeds, five communication rounds, and clean
and sign-flip conditions. Our results are mixed: contribution-aware aggregation improves
paired PR-AUC on MLG-ULB but not on the 10,000-row Kaggle dataset. The
heterogeneous encoder comparison is uncertain on MLG-ULB and favors
encoder+FedAvg on the smaller Kaggle dataset. With five seeds and only 23
fraud cases in the smaller test split, we treat these findings as exploratory,
not evidence of general superiority. We simulate clients as feature subsets,
and our prototype is server-coordinated; it does not provide secure
aggregation or differential privacy.
\end{abstract}

\begin{keywords}
Federated learning, fraud detection, heterogeneous features, robust
aggregation, client contribution, imbalanced classification
\end{keywords}

\titlepgskip=-15pt

\maketitle

\section{Introduction}
Fraud is rare in transaction data, so accuracy alone can obscure how well a
model finds it; precision--recall analysis is more informative in this
setting \cite{saito2015}. At the same time, financial institutions may not be
able to pool their transaction records. Federated learning offers a way to
train a shared model while keeping raw records at their source
\cite{mcmahan2017}.

Standard federated averaging (FedAvg) assumes that participating clients can
train compatible parameterizations \cite{mcmahan2017}. Statistical
heterogeneity can cause client drift, while feature-schema heterogeneity makes
direct averaging of complete model states ill-defined. Robust aggregation
methods such as coordinate-wise median, trimmed mean, and Krum have also been
studied under adversarial or Byzantine update assumptions
\cite{yin2018,blanchard2017}. These methods address different failure modes
and should not be interpreted as universal defenses.

In our setting, the obstacle is not only the distribution of client data: clients may
also have different input features, so their full model parameters cannot be
averaged directly. We keep a feature-specific encoder at each client and
aggregate a shared prediction torso. We then compare standard aggregation
rules with a contribution score that combines local performance and update
reliability with estimated usefulness to other clients. This is an empirical
test of a heuristic, not a convergence or privacy guarantee.

We evaluate the current implementation using newly generated benchmark
artifacts; we do not include results from earlier project iterations.

\subsection{Study Scope}
We focus on three practical questions:
\begin{itemize}
  \item Can clients with different feature subsets train through local
  encoders and a common prediction torso?
  \item How does adaptive paired scoring compare with encoder-based FedAvg
  and same-schema aggregation baselines?
  \item Do the comparisons change under a simulated sign-flip update?
\end{itemize}

\section{Related Work}
\subsection{Federated Averaging and Heterogeneous Optimization}
FedAvg averages locally trained model parameters and can reduce the number of
communication rounds relative to synchronized stochastic gradient descent
\cite{mcmahan2017}. FedProx adds a proximal term to local training to address
heterogeneity \cite{li2020}. SCAFFOLD uses control variates to reduce client
drift \cite{karimireddy2020}; FedNova normalizes local updates to address
objective inconsistency when clients perform unequal amounts of work
\cite{wang2020}. Only FedAvg and FedProx are included in the same-schema
comparison here; SCAFFOLD and FedNova provide context, not measured baselines.

\subsection{Robust Aggregation}
Krum selects an update using distances to neighboring updates
\cite{blanchard2017}. Coordinate median and trimmed-mean rules have been
analyzed for Byzantine-robust distributed learning \cite{yin2018}. We include
trimmed mean and one synthetic sign-flip condition. That experiment is narrow:
it does not establish resistance to adaptive or targeted poisoning.

\subsection{Heterogeneous Federated Models and Client Utility}
The encoder/shared-torso design gives clients with different tabular schemas
a common set of parameters to aggregate. Feature-keyed initialization and a
local proximal anchor are implementation choices; neither shows that the
clients' latent representations are aligned. Likewise, counterfactual
validation utility is only a proxy for cross-client usefulness. It is not a
causal estimate of benefit to another institution. Data Shapley is one
principled approach to data valuation, but the heuristic evaluated here is
not a Shapley estimator and does not provide its axiomatic guarantees
\cite{ghorbani2019}.

\section{Method}
\subsection{Problem Setup}
Let client $k$ hold local rows $D_k$ with feature set $\mathcal{F}_k$ and binary
label $y$. Each client has a local encoder $e_k(\cdot;\phi_k)$ and shares a
torso $g(\cdot;\theta)$:
\begin{equation}
  \hat{y}=g(e_k(x_{\mathcal{F}_k};\phi_k);\theta).
\end{equation}
The encoder parameters $\phi_k$ remain client-specific. At each communication
round, clients train locally and transmit torso updates; the coordinator
aggregates only compatible $\theta$ tensors. Fig.~\ref{fig:architecture}
summarizes the data and parameter flow.
For a standard sample-weighted round, let $n_k=|D_k|$ and let
$\theta_{k}^{(t+1)}$ denote client $k$'s returned torso. The server update is
\begin{equation}
  \theta^{(t+1)}=\sum_{k=1}^{K}\frac{n_k}{\sum_{j=1}^{K}n_j}
  \theta_k^{(t+1)}.
\end{equation}
Contribution-aware runs replace the sample proportions with scorer-derived
weights while still averaging only the shared torso.

\begin{figure}[t]
\centering
\benchmarkfigure{architecture.pdf}{width=\columnwidth}
\caption{\architecturecaption{Author-created schematic. Local encoders remain
at their simulated clients. The coordinator aggregates the shared torso using
client weights from the contribution score; the score uses local evidence
and a common reference split.}}
\label{fig:architecture}
\end{figure}

\subsection{Adaptive Paired Contribution Score}
The scorer compares client updates using validation information and a
coordinator-held reference split. Its local pair combines validation quality
with update-scale reliability:
\begin{equation}
  L_k=\frac{Q_k+T_k}{2},
\end{equation}
where $Q_k$ is the robustly standardized client validation PR-AUC, passed
through a logistic transform. To reduce the influence of a small number of
validation fraud cases, a client's PR-AUC is shrunk toward the participating
clients' median before standardization:
\begin{equation}
  \widetilde{p}_k=\operatorname{median}(p)
  +\frac{m_k}{m_k+20}\bigl(p_k-\operatorname{median}(p)\bigr),
\end{equation}
where $p_k$ and $m_k$ are its validation PR-AUC and validation fraud count.
The reliability term $T_k$ summarizes whether the torso-update norm is an
extreme outlier relative to other clients in the same round. With
$z_k=\log(\lVert\theta_k-\theta^{(t)}\rVert_2+\epsilon)$, the implementation
uses the median $c$ of the $z_k$ values and robust scale
$s=\max(1.4826\,\operatorname{MAD}(z),0.35)$:
\begin{align}
  a_k &= \max\bigl(|z_k-c|-2s,0\bigr),\\
  T_k^{\mathrm{round}} &= \exp(-a_k/s),\\
  T_k &= 0.75T_k^{\mathrm{round}}+0.25T_k^{\mathrm{prior}}.
\end{align}
This is a relative update-scale heuristic, not a test of whether an update is
malicious.

The utility pair is
\begin{equation}
  U_k=\frac{N_k+C_k}{2},
\end{equation}
where $N_k$ measures improvement on hard positive reference examples and
$C_k$ measures leave-one-client-out reference PR-AUC gain. The full candidate
state is a sample-weighted average of the received updates; for each client,
the counterfactual state is the corresponding average with that client's
update omitted. Complementarity is the full-state minus counterfactual
average precision:
\begin{equation}
  \Delta_k^{\mathrm{AP}}
  =\operatorname{AP}(\theta_{\mathrm{all}})
  -\operatorname{AP}(\theta_{\mathrm{without}\ k}).
\end{equation}
The hard-positive set $\mathcal{H}$ contains reference fraud cases assigned
the lowest probabilities by the current global model, up to the larger of
eight cases or half the available positives. The novelty gain is
\begin{equation}
  \Delta_k^{\mathrm{hard}}
  =\frac{1}{|\mathcal{H}|}\sum_{i\in\mathcal{H}}
  \log p_{\mathrm{all},i}
  -\frac{1}{|\mathcal{H}|}\sum_{i\in\mathcal{H}}
  \log p_{\mathrm{without}\ k,i}.
\end{equation}
For each utility gain, the score uses the 10th-percentile lower bound from
paired class-stratified bootstrap resamples. That lower bound is divided by
the median absolute client lower bound for the round and passed through a
logistic transform, giving normalized $N_k$ and $C_k$ values in $[0,1]$.
Only eight bootstrap resamples were used per estimate in this pilot, so these
lower bounds are coarse and should not be read as precise uncertainty
intervals.

The temporal signal $H_k$ is initialized at $0.5$. Thereafter, the raw
complementarity lower bound is tracked with a $0.7$ exponential-moving-average
decay and transformed by a logistic function with scale $0.01$. For the
adaptive-pair profile, the client-specific gate is
\begin{align}
  s_k &= 0.7C_k+0.3H_k,\\
  \lambda_k &= \operatorname{clip}(0.2+0.6s_k,0.2,0.8),\\
  R_k &= (1-\lambda_k)L_k+\lambda_kU_k.
\end{align}
The client aggregation weights are
\begin{equation}
  w_k=\frac{\exp(R_k/\tau)}{\sum_j\exp(R_j/\tau)},\qquad \tau=0.35.
\end{equation}
The softmax makes the weights nonnegative and sum to one; $\tau$ controls how
strongly score differences affect the final weights. These weights are
recomputed for each round.
The gate is heuristic and can be affected by reference-set composition,
calibration, and small client validation samples. The same-schema
contribution-aware baseline uses the five scorer dimensions with fixed
coefficients (0.20, 0.20, 0.20, 0.25, and 0.15 for quality, trust, novelty,
complementarity, and temporal utility, respectively); it is distinct from the
adaptive-pair profile. In particular, results for that fixed-weight baseline
must not be attributed to the adaptive gate used in the heterogeneous
comparison.

\subsection{Attack and Baseline Conditions}
The same-schema experiment compares FedAvg, FedProx, coordinate-wise trimmed
mean, and five-signal contribution-aware aggregation. The heterogeneous
experiment compares encoder+FedAvg, encoder+adaptive paired scoring, and a
full-schema FedAvg contextual reference. The latter uses more features and is
not a matched-schema baseline.

In the sign-flip condition, client 0's model delta is negated around the
current global state at each round. The clean condition applies no update
transformation. We use this synthetic perturbation to isolate a limited behavior; it does not
model a complete adversary.

\section{Experimental Design}
\subsection{Datasets and Preprocessing}
We downloaded the Kaggle ``Credit Card Fraud Detection'' collection
associated with MLG-ULB \cite{creditcarddataset}. Our downloaded copy contains
284,807 rows, 492 fraud labels, and 30 numeric predictors. We also downloaded
the Kaggle ``Credit Card Fraud Detection Dataset'' \cite{miadul_dataset}. The
copy we used contains 10,000 rows, 151 fraud labels, and a different set of
transaction features. We excluded its transaction identifier, one-hot encoded
its merchant category, and fit numeric transformations on the training
partition only. We evaluate the datasets separately and do not pool their
metrics.

Table~\ref{tab:datasets} summarizes the raw sample sizes and model input
dimensions. Table~\ref{tab:bas1-design} records the complete comparison
structure and run protocol for this expanded presentation.
\begin{table}[t]
\caption{Characteristics of the Kaggle datasets used in our experiments}
\label{tab:datasets}
\centering
\begin{tabular}{lrrrr}
\toprule
Dataset & Rows & Fraud & Rate & Inputs \\
\midrule
MLG-ULB (Kaggle) & 284,807 & 492 & 0.173\% & 30 \\
Kaggle 10,000-row & 10,000 & 151 & 1.510\% & 12 \\
\bottomrule
\end{tabular}
\end{table}

\begin{table*}[t]
\caption{Benchmark design and run protocol}
\label{tab:bas1-design}
\centering
\begin{tabular}{p{0.20\textwidth}p{0.73\textwidth}}
\toprule
Aspect & Implemented setting \\
\midrule
Same-schema track & FedAvg, FedProx, coordinate-wise trimmed mean, and fixed-weight contribution-aware aggregation. \\
Heterogeneous track & Client-specific encoders with shared-torso FedAvg or adaptive paired scoring; full-schema FedAvg is an unmatched contextual reference. \\
Client simulation & Five clients; label-skewed Dirichlet partitions ($\alpha=0.1$) for the standard track and feature-subset clients for the heterogeneous track. \\
Training schedule & Seeds 42--46, five communication rounds, one local epoch per round, and batch size 1024. \\
Conditions & Clean training and a synthetic sign-flip that negates client 0's model delta at each round. \\
Scorer uncertainty & Eight class-stratified bootstrap draws per utility estimate; coarse pilot estimates only. \\
Evaluation & Fixed test split per dataset; PR-AUC is primary and ROC-AUC secondary. Intervals summarize five seeds, not test-sample uncertainty. \\
\bottomrule
\end{tabular}
\end{table*}


We stratified the 10,000-row Kaggle dataset into 70\% training, 15\%
validation, and 15\% test partitions. For MLG-ULB, we used its existing
stratified partitions. We held both test partitions out until final evaluation.
We ran five seeds and five communication rounds, with five simulated clients
and one local epoch per round. Standard-track clients use label-skewed
Dirichlet partitions ($\alpha=0.1$); the heterogeneous track uses five
feature-subset clients with local encoders. We used batch size 1024, two CPU
threads, and an NVIDIA GeForce RTX 5050 Laptop GPU. FedProx uses $\mu=0.01$,
and the heterogeneous encoder alignment penalty uses
$\mu_{\mathrm{align}}=0.01$. Our contribution scorer uses eight
class-stratified bootstrap draws per utility estimate in this pilot. We retain
the software, split hashes, per-seed configuration, and run status in the
result artifacts.

\subsection{Metrics and Statistical Reporting}
PR-AUC is the primary ranking metric because the positive class is rare and
precision--recall analysis is informative under strong class imbalance
\cite{saito2015}.
We compute average precision from the precision--recall curve:
\begin{equation}
  \mathrm{AP}=\sum_n(R_n-R_{n-1})P_n,
\end{equation}
where $P_n$ and $R_n$ are precision and recall at threshold $n$. This
step-wise average-precision definition is used consistently for client
validation and final test evaluation.
Both experiment tracks report PR-AUC and ROC-AUC. The same-schema track also
reports threshold-dependent F1, precision, and recall; these threshold metrics
are not present in the heterogeneous-run artifacts.
We report means and two-sided $t$ intervals across seeds. We form paired
method deltas using the same seed values. These intervals describe
training/partition variability on the fixed test rows; they do not quantify
test-sample uncertainty. No multiple-comparison correction is applied.
For five seed-level test scores $x_s$, the reported mean and interval are
\begin{align}
  \bar{x} &= \frac{1}{5}\sum_{s=1}^{5}x_s,\\
  \mathrm{CI}_{95\%} &=
  \bar{x}\ \pm\ t_{0.975,4}\frac{s_x}{\sqrt{5}},
\end{align}
where $s_x$ is the sample standard deviation across seeds. For each matched
comparison, the seed-wise difference is formed before summarizing:
\begin{equation}
  \Delta_s = x_{s,\mathrm{candidate}}-x_{s,\mathrm{reference}}.
\end{equation}
Thus the paired interval reflects variation in five matched training runs on
the same fixed test rows; it does not measure uncertainty from drawing a new
test sample.

\section{Results}
We completed all eight planned benchmark jobs. Figs.~\ref{fig:clean-pr-ranks}
and \ref{fig:signflip-pr-ranks} rank methods by mean test PR-AUC, with a
separate figure for each attack condition. Within each dataset panel, ranks
are calculated separately for the same-schema and heterogeneous tracks; the
two tracks are not pooled into one leaderboard. The paired plots in
Figs.~\ref{fig:clean-paired} and \ref{fig:signflip-paired}, together with
Table~\ref{tab:paired-results}, compare contribution-aware methods against
their matched FedAvg references using the same seeds. Table~\ref{tab:bas1-mean}
adds the per-method PR-AUC means and seed intervals. Fig.~\ref{fig:bas1-seed-scatter}
shows the matched seed-level scores against the identity line. ROC-AUC is shown
as a secondary metric in separate attack-specific plots.

\begin{figure*}[t]
\centering
\benchmarkfigure{bas1_ranked_none_pr_auc.pdf}{width=0.93\textwidth}
\caption{\resultcaption{Clean-condition test PR-AUC rankings for both datasets. In each
dataset, ranks are ordered by the five-seed mean within the same-schema and
heterogeneous tracks separately. Bars show mean PR-AUC and whiskers show
two-sided 95\% $t$ intervals over seeds 42--46 on the fixed test split.
Full-schema FedAvg is contextual only because it uses more features. Rank
order is descriptive and is not a significance test. Source: our
visualization of completed benchmark artifacts; intervals are clipped to
[0, 1] for display. Interpretation: on MLG-ULB, contribution-aware
aggregation (0.719) and trimmed mean (0.718) are nearly tied in the
same-schema track. On the 10,000-row data, FedAvg and FedProx are near 0.391,
while contribution-aware aggregation is 0.281.}}
\label{fig:clean-pr-ranks}
\end{figure*}

In the clean MLG-ULB panel, contribution-aware aggregation has the highest
same-schema mean PR-AUC (0.719), but trimmed mean is nearly tied (0.718).
Within the heterogeneous track, adaptive pairs reach 0.702, compared with
0.687 for encoder+FedAvg; the full-schema model is an unmatched reference.
The 10,000-row dataset gives a different ranking: FedAvg and FedProx are
near 0.391 in the same-schema track, while contribution-aware aggregation
reaches 0.281. In its heterogeneous track, the adaptive-pair mean is 0.265,
below encoder+FedAvg at 0.302. Thus the clean ranking is dataset-dependent;
the bars describe five-seed means, not statistically established superiority.

\begin{figure*}[t]
\centering
\benchmarkfigure{bas1_ranked_sign_flip_pr_auc.pdf}{width=0.93\textwidth}
\caption{\resultcaption{Test PR-AUC rankings under the simulated sign-flip condition, shown
separately from clean training. Within each dataset, methods are ranked by
five-seed mean within their respective track. Whiskers are two-sided 95\%
$t$ intervals over seeds 42--46 on the fixed test split. Full-schema FedAvg
is contextual only. Rank order is descriptive. Source: our
visualization of completed benchmark artifacts; intervals are clipped to
[0, 1] for display. Interpretation: the
MLG-ULB same-schema means for contribution-aware aggregation and trimmed mean
remain close (0.715 and 0.714). On the 10,000-row data, FedAvg averages 0.318
and contribution-aware aggregation 0.219.}}
\label{fig:signflip-pr-ranks}
\end{figure*}

Under sign flip, contribution-aware aggregation and trimmed mean remain
closely ranked on MLG-ULB (mean PR-AUC 0.715 and 0.714). The heterogeneous
adaptive-pair result is 0.705, while encoder+FedAvg is much more variable
across seeds; its wide interval makes the apparent gap uncertain. On the
10,000-row dataset, same-schema FedAvg averages 0.318 versus 0.219 for
contribution-aware aggregation. In the heterogeneous track, encoder+FedAvg
averages 0.253 and adaptive pairs 0.212. These results do not establish
resistance to poisoning: the tested sign flip affects one simulated client
and is only a narrow perturbation.

\begin{figure*}[t]
\centering
\benchmarkfigure{bas1_ranked_none_roc_auc.pdf}{width=0.93\textwidth}
\caption{\resultcaption{Clean-condition test ROC-AUC rankings for both datasets. Within each
dataset, ranks are calculated separately for the same-schema and
heterogeneous tracks. Bars show five-seed means and whiskers show two-sided
95\% $t$ intervals on the fixed test split. Full-schema FedAvg is a
contextual, unmatched reference. Source: our visualization of completed
benchmark artifacts;
intervals are clipped to [0, 1] for display. Interpretation: MLG-ULB
same-schema ROC-AUC means are tightly clustered (0.954--0.958). On the
10,000-row data, FedAvg and FedProx are about 0.906, versus 0.868 for
contribution-aware aggregation.}}
\label{fig:clean-roc-ranks}
\end{figure*}

ROC-AUC provides a complementary, less class-imbalance-sensitive ranking.
On MLG-ULB, same-schema methods cluster between 0.954 and 0.958, so their
ordering differs little in practical terms; the 10,000-row same-schema
FedAvg/FedProx means are about 0.906, versus 0.868 for contribution-aware
aggregation. In the smaller dataset's heterogeneous track, adaptive pairs
have a mean of 0.911 and encoder+FedAvg 0.902. This differs from their
PR-AUC ordering, underscoring why ROC-AUC alone is insufficient for a rare
fraud class. Intervals still reflect five seeds on fixed test rows.

\begin{figure*}[t]
\centering
\benchmarkfigure{bas1_ranked_sign_flip_roc_auc.pdf}{width=0.93\textwidth}
\caption{\resultcaption{Test ROC-AUC rankings under the simulated sign-flip condition.
Within each dataset, ranks are calculated separately for the same-schema and
heterogeneous tracks. Bars show five-seed means and whiskers show two-sided
95\% $t$ intervals on the fixed test split. Full-schema FedAvg is a
contextual, unmatched reference. Source: our visualization of completed
benchmark artifacts;
intervals are clipped to [0, 1] for display. Interpretation: on MLG-ULB,
FedAvg and FedProx are near 0.959 ROC-AUC, while contribution-aware
aggregation is 0.954. On the 10,000-row data, same-schema FedAvg averages
0.864 and contribution-aware aggregation 0.816; adaptive pairs lead
encoder+FedAvg in the heterogeneous track (0.881 versus 0.851).}}
\label{fig:signflip-roc-ranks}
\end{figure*}

The sign-flip ROC-AUC rankings are similarly mixed. On MLG-ULB, FedAvg and
FedProx are near 0.959, while contribution-aware aggregation is 0.954. On
the 10,000-row data, same-schema FedAvg averages 0.864 and contribution-aware
aggregation 0.816; trimmed mean falls to 0.726. The heterogeneous
adaptive-pair mean is 0.881, compared with 0.851 for encoder+FedAvg. These
are descriptive means: large intervals on several methods, especially the
full-schema MLG-ULB reference, preclude a strong ranking claim.

\begin{figure*}[!tp]
\centering
\benchmarkfigure{bas1_paired_none_pr_auc.pdf}{width=0.93\textwidth}
\caption{\resultcaption{Paired change in clean-condition test PR-AUC for each
contribution-aware method relative to its matched FedAvg reference. Positive
values favor the candidate; the dashed line marks no change. Small points
show per-seed paired differences; diamonds and whiskers show the mean and
two-sided 95\% $t$ interval. Source: our analysis of completed benchmark
artifacts, seeds
42--46; test rows are fixed within each dataset. Interpretation: the
same-schema MLG-ULB change is positive, while both estimates for
the 10,000-row Kaggle dataset are negative. The heterogeneous MLG-ULB interval is very wide and crosses
zero, so our results do not show a consistent gain under this perturbation.}}
\label{fig:clean-paired}
\end{figure*}

The paired plot isolates the change from each matched FedAvg baseline,
removing between-seed level differences. Positive values on MLG-ULB favor
contribution-aware aggregation in the same-schema track; the heterogeneous
estimate is smaller and its interval crosses zero. Both paired estimates for
the 10,000-row Kaggle dataset are negative; the clean heterogeneous interval
excludes zero. This view makes the cross-dataset reversal clearer than rank alone,
while remaining an exploratory five-seed comparison.

\begin{figure*}[!tp]
\centering
\benchmarkfigure{bas1_paired_sign_flip_pr_auc.pdf}{width=0.93\textwidth}
\caption{\resultcaption{Paired change in sign-flip-condition test PR-AUC for each
contribution-aware method relative to its matched FedAvg reference. Positive
values favor the candidate; the dashed line marks no change. Small points
show per-seed paired differences; diamonds and whiskers show the mean and
two-sided 95\% $t$ interval. Source: our analysis of completed benchmark
artifacts, seeds
42--46; test rows are fixed within each dataset.}}
\label{fig:signflip-paired}
\end{figure*}

\begin{figure*}[t]
\centering
\benchmarkfigure{bas1_seed_scatter_pr_auc.pdf}{width=0.93\textwidth}
\caption{\resultcaption{Seed-paired PR-AUC scores for each contribution-aware
method against its matched FedAvg reference. Each point is one of seeds
42--46; the dashed identity line indicates equal candidate and reference
scores. Panels separate datasets, colors distinguish comparison tracks, and
marker shapes distinguish clean and sign-flip conditions. These five points
are descriptive and are not independent test samples.}}
\label{fig:bas1-seed-scatter}
\end{figure*}


\begin{table*}[t]
\caption{Paired change in test PR-AUC relative to the matched FedAvg baseline
(encoder+FedAvg for the heterogeneous track). Source: per-seed benchmark
artifacts, matched by random seed.}
\label{tab:paired-results}
\centering
\begin{tabular}{llrr}
\toprule
Dataset & Condition & Same-schema $\Delta$PR-AUC (95\% CI) &
Heterogeneous $\Delta$PR-AUC (95\% CI) \\
\midrule
MLG-ULB & Clean & $+0.040$ ($[0.013,\,0.068]$) & $+0.015$ ($[-0.012,\,0.042]$) \\
MLG-ULB & Sign-flip & $+0.036$ ($[0.014,\,0.058]$) & $+0.260$ ($[-0.160,\,0.680]$) \\
Kaggle 10,000-row & Clean & $-0.110$ ($[-0.258,\,0.038]$) & $-0.037$ ($[-0.069,\,-0.004]$) \\
Kaggle 10,000-row & Sign-flip & $-0.098$ ($[-0.190,\,-0.007]$) & $-0.041$ ($[-0.133,\,0.052]$) \\
\bottomrule
\end{tabular}
\end{table*}

\begin{table*}[t]
\caption{Mean test PR-AUC and two-sided 95\% $t$ intervals across five seeds. The same-schema and heterogeneous columns are separate tracks; the full-schema contextual baseline is omitted.}
\label{tab:bas1-mean}
\centering
\begin{tabular}{llcccc}
\toprule
Dataset & Condition & FedAvg & Contribution-aware & Encoder+FedAvg & Encoder+adaptive pairs \\
\midrule
MLG-ULB & Clean & 0.678 $\pm$ 0.026 & 0.719 $\pm$ 0.003 & 0.687 $\pm$ 0.032 & 0.702 $\pm$ 0.017 \\
MLG-ULB & Sign-flip & 0.679 $\pm$ 0.019 & 0.715 $\pm$ 0.004 & 0.445 $\pm$ 0.419 & 0.705 $\pm$ 0.012 \\
Kaggle 10,000-row & Clean & 0.391 $\pm$ 0.130 & 0.281 $\pm$ 0.207 & 0.302 $\pm$ 0.189 & 0.265 $\pm$ 0.176 \\
Kaggle 10,000-row & Sign-flip & 0.318 $\pm$ 0.196 & 0.219 $\pm$ 0.182 & 0.253 $\pm$ 0.224 & 0.212 $\pm$ 0.143 \\
\bottomrule
\end{tabular}
\end{table*}


\subsection{Clean Condition}
On the MLG-ULB dataset, the same-schema contribution-aware method had a mean
paired PR-AUC increase of 0.040 under clean training and 0.036 under
sign-flip updates. Both five-seed intervals were above zero. The heterogeneous
adaptive-pair method had a smaller clean-data estimate of 0.015; its interval
included zero. Under sign flip its estimate was 0.260, but the interval was
wide and also included zero, reflecting substantial seed-to-seed variation.
By mean PR-AUC, contribution-aware aggregation ranked first in both
same-schema conditions; trimmed mean ranked a close second. In the clean
heterogeneous panel, adaptive pairs ranked first, while under sign flip the
same method again ranked first, alongside a much wider interval for the
encoder+FedAvg reference.

The 10,000-row Kaggle dataset shows a different pattern. Relative to
same-schema FedAvg, contribution-aware aggregation had a paired PR-AUC change
of -0.110 under clean training and -0.098 under sign flip. For the
heterogeneous track, adaptive pairs were below encoder+FedAvg by 0.037 in the
clean condition and 0.041 under sign flip; only the clean interval excluded
zero. We therefore find no consistent cross-dataset improvement from
contribution-aware weighting in this pilot. FedProx and FedAvg had the highest
same-schema means in the clean and sign-flip panels, respectively.
Encoder+FedAvg ranked above adaptive pairs in both heterogeneous panels. These
are rankings of observed means, not evidence that closely spaced methods
differ reliably.

\subsection{Sign-Flip Condition}
We observed smaller ROC-AUC changes, with no consistent direction across
datasets and method tracks. On the 10,000-row dataset, just 23 fraud examples
fall in the held-out test partition. Its PR-AUC intervals are consequently
wide, and small changes in the ranking of a few positive examples can have a
visible effect. The synthetic sign-flip condition should not be read as a
general test of poisoning resistance.

\subsection{Encoder Compatibility on the 10,000-Row Kaggle Dataset}
We completed the encoder experiments on both datasets. On the 10,000-row
Kaggle dataset, encoder+FedAvg exceeded adaptive pairs in mean PR-AUC in both
conditions.
Full-schema FedAvg appears as a contextual reference in
Figs.~\ref{fig:clean-pr-ranks} and \ref{fig:signflip-pr-ranks}; it sees all
input features and is not a matched-schema comparison.

\section{Discussion}
We evaluated both input schemas, but our results do not point to one best
aggregation rule across them. Contribution-aware aggregation improved PR-AUC
over same-schema FedAvg on the MLG-ULB split; the paired differences went the
other way on the 10,000-row Kaggle dataset. The heterogeneous comparison was
less decisive on MLG-ULB and favored encoder+FedAvg on the smaller dataset.
These observations come from five seeds and one fixed test split per dataset;
they do not show that one method will generalize to other banks or time
periods.

The 10,000-row Kaggle dataset has only 23 fraud examples in its held-out test
partition. Its PR-AUC and threshold metrics are therefore sensitive to a small
number of examples. Our fixed split and five seeds do not replace repeated
stratified test sampling, temporal evaluation, or external validation.

\section{Threats to Validity and Limitations}
\begin{itemize}
  \item The simulated clients are feature-subset partitions of pooled
  datasets, not separate real financial institutions.
  \item The coordinator architecture is not peer-to-peer decentralized FL.
  \item The experiment does not provide differential privacy, secure
  aggregation, encryption, or a formal privacy guarantee.
  \item Feature-keyed initialization and proximal anchoring do not validate
  latent-space alignment.
  \item Contribution scoring queries a common reference set in the simulator;
  it is not a deployable privacy-preserving client-valuation protocol.
  \item Five seeds, five rounds, one fixed test split per dataset, and one
  simple attack condition are exploratory and insufficient for broad
  superiority or security claims.
  \item The paired comparisons are exploratory and no adjustment for multiple
  method comparisons was applied.
  \item We evaluate selected comparators; we do not implement SCAFFOLD or
  FedNova in these experiments.
\end{itemize}

\section{Conclusion}
Our prototype combines client-specific tabular encoders, a shared torso, and
adaptive paired contribution scoring. We conduct a limited comparison on two
fixed datasets, simulate clients as feature subsets, and test one synthetic
sign-flip condition. We limit our performance conclusions to these settings.
A stronger evaluation would require independently held data, temporal
testing, more seeds, realistic client ownership, and evidence that the learned
client representations are aligned.

\section*{Funding}
We received no external funding for this research.

\section*{Data Availability}
We downloaded both datasets used in this study from their respective Kaggle
listings: the MLG-ULB credit-card dataset [9] and the 10,000-row ``Credit Card
Fraud Detection Dataset'' [10]. We used local copies of those downloads.
Readers can access the dataset listings through the cited links; any
redistribution of dataset files remains subject to the terms and licenses on
their respective Kaggle pages. We will share processed data and run artifacts
only when their licenses and privacy conditions permit.

\clearpage
\begin{thebibliography}{00}
\bibitem{saito2015}
T. Saito and M. Rehmsmeier, ``The precision-recall plot is more informative
than the ROC plot when evaluating binary classifiers on imbalanced datasets,''
\emph{PLOS ONE}, vol. 10, no. 3, Art. no. e0118432, Mar. 2015,
\href{https://doi.org/10.1371/journal.pone.0118432}
{doi: 10.1371/journal.pone.0118432}.

\bibitem{mcmahan2017}
B. McMahan, E. Moore, D. Ramage, S. Hampson, and B. A. y Arcas,
``Communication-efficient learning of deep networks from decentralized
data,'' in \emph{Proc. 20th Int. Conf. Artif. Intell. Stat.}, PMLR, vol. 54,
2017, pp. 1273--1282.
\href{https://proceedings.mlr.press/v54/mcmahan17a.html}
{[Online]. Available: PMLR.}

\bibitem{yin2018}
D. Yin, Y. Chen, R. Kannan, and P. Bartlett, ``Byzantine-robust distributed
learning: Towards optimal statistical rates,'' in \emph{Proc. ICML}, 2018,
PMLR, vol. 80, pp. 5650--5659.
\href{https://proceedings.mlr.press/v80/yin18a.html}
{[Online]. Available: PMLR.}

\bibitem{blanchard2017}
P. Blanchard, E. M. El Mhamdi, R. Guerraoui, and J. Stainer,
``Machine learning with adversaries: Byzantine tolerant gradient descent,''
in \emph{Proc. NeurIPS}, 2017.
\href{https://proceedings.neurips.cc/paper/2017/hash/f4b9ec30ad9f68f89b29639786cb62ef-Abstract.html}
{[Online]. Available: NeurIPS.}

\bibitem{li2020}
T. Li, A. K. Sahu, M. Zaheer, M. Sanjabi, A. Talwalkar, and V. Smith,
``Federated optimization in heterogeneous networks,'' in
\emph{Proc. 3rd MLSys Conf.}, 2020.
\href{https://proceedings.mlsys.org/paper_files/paper/2020/hash/1f5fe83998a09396ebe6477d9475ba0c-Abstract.html}
{[Online]. Available: MLSys.}

\bibitem{karimireddy2020}
S. P. Karimireddy, S. Kale, M. Mohri, S. Reddi, S. Stich, and A. T. Suresh,
``SCAFFOLD: Stochastic controlled averaging for federated learning,'' in
\emph{Proc. 37th Int. Conf. Mach. Learn.}, PMLR, vol. 119, 2020,
pp. 5132--5143.
\href{https://proceedings.mlr.press/v119/karimireddy20a.html}
{[Online]. Available: PMLR.}

\bibitem{wang2020}
J. Wang, Q. Liu, H. Liang, G. Joshi, and H. V. Poor, ``Tackling the
objective inconsistency problem in heterogeneous federated optimization,''
in \emph{Proc. Adv. Neural Inf. Process. Syst.}, vol. 33, 2020.
\href{https://proceedings.neurips.cc/paper/2020/hash/564127c03caab942e503ee6f810f54fd-Abstract.html}
{[Online]. Available: NeurIPS.}

\bibitem{ghorbani2019}
A. Ghorbani and J. Zou, ``Data Shapley: Equitable valuation of data for
machine learning,'' in \emph{Proc. 36th Int. Conf. Mach. Learn.}, PMLR,
vol. 97, 2019, pp. 2242--2251.
\href{https://proceedings.mlr.press/v97/ghorbani19c.html}
{[Online]. Available: PMLR.}

\bibitem{creditcarddataset}
MLG-ULB, ``Credit Card Fraud Detection,'' Kaggle dataset.
Available: \href{https://www.kaggle.com/datasets/mlg-ulb/creditcardfraud}
{Kaggle dataset page}.
Accessed: Nov. 6, 2026.

\bibitem{miadul_dataset}
Miadul, ``Credit Card Fraud Detection Dataset,'' Kaggle.
Available: \href{https://www.kaggle.com/datasets/miadul/credit-card-fraud-detection-dataset/data}
{Kaggle dataset page}.
Accessed: Nov. 6, 2026.
\end{thebibliography}

\EOD
\end{document}
